% ===== 原书 p.1 =====
\section{引言}

考察高维流形的典型方法之一，是考虑纤维化。具有严格意义下的纤维丛结构的流形并不多，但若允许奇异纤维，适用范围便会急剧扩大。

在辛流形上引入纤维化结构时，自然要对纤维也附加某种辛几何条件。辛流形的辛几何子流形中，首先能想到的是辛子流形与 Lagrange 子流形。

辛流形以辛子流形为纤维的纤维化中，最重要的是辛 Lefschetz 纤维化。它是只允许非退化奇点的纤维丛的推广，也是 Morse 函数在伪全纯几何中的类比。这里，所谓从殆复流形到 Riemann 面的伪全纯映射非退化，是指适当地取殆复局部坐标后可以写成 $(z_1,\dots,z_n)\mapsto z_1^2+\cdots+z_n^2$ 的形式。此外，我们预先在辛流形上取定一个与辛结构相容的殆复结构。由 Donaldson 的著名定理，任意辛流形都具有辛 Lefschetz 束（symplectic Lefschetz pencil）的结构，并且经过适当的爆破，可以赋予辛 Lefschetz 纤维化的结构。
% [Tutor Note: 「適当に爆発することで」指沿 pencil 的基点集作爆破；Lefschetz 束与 Lefschetz 纤维化的差别正在于前者容许基点，爆破掉基点即得纤维化]

辛流形的辛子流形，由辛形式的限制非退化来定义；而 Lagrange 子流形的概念在某种意义上恰好与之相对，由辛形式的限制为零来定义。非退化性是开条件，而为零是闭条件，因此一个子流形是 Lagrange 的，比它是辛的严格得多。辛子流形经过少许变形仍是辛子流形，Lagrange 子流形经过少许变形却不再是 Lagrange 子流形。

在由 Lagrange 子流形给出的纤维化中，一般纤维为环面的情形尤为重要。其根据之一是 Liouville–Arnold 定理，它断言 Liouville 意义下的完全可积系统在适当条件下给出 Lagrange 环面纤维化。另一个根据是，辛流形的几何量子化所得的 Hilbert 空间，与以 Lagrange 环面纤维化的 Bohr--Sommerfeld 纤维为基的向量空间密切相关。这一关系在环面流形上最为显著，但在旗流形以及 Riemann 面上秩 $2$ 的向量丛的模空间上，也能发现类似的结构。
% [Tutor Note: 几何量子化的「极化」选取决定量子化 Hilbert 空间；Bohr--Sommerfeld 条件把可容许的极化几何化为某个格的条件，这正是纤维化与量子化相接之处]

下面介绍这种 Lagrange 环面纤维化与镜像对称的深刻关联。
% [Tutor Note: 原文此处印作「ファーブレーション」，是「ファイブレーション」的误植，译文按正字处理]

% CHUNK c00 DONE pages p1
